%! Author = lili
%! Date = 5/6/26

\section{Nukleosomen und Histonvarianten}\label{sec:nukleosomen-und-histonvarianten}

% begriffe für Karteikarten: Histonoktamer,
% Ich sollte laut reden, nicht ablesen und nicht gestelzt Begriffe wie Histonoktamer statt Histon benutzen
% Er fragt gerne, was das mit der Bioinformatik zu tun hat

%\paragraph{Sollten wir die Histonvarianten für die Klausur kennen?}

\paragraph{F 7: Welche Variante macht was oder macht jede Variante alles davon?}  Nicht jede macht alles, kann man
auf anderer Folie sehen, war da nur nicht mit drauf was was ist.

\paragraph{F 8: was meinst du mit Histone "markieren" verschiedene Transkriptionszustände?} Man könnte es einfach
daran erkennen.
Theoretisch, aber auch man könnte auch Antikörper produzieren lassen, die an bestimmte
Modifikationen binden und sie so praktisch markieren.

\paragraph{F 8: Ich habe nicht verstanden was ein +1 Nukleosom ist. Ist das da abgebildet?} Ja, ganz klein.
Das ist einfach das, was an der +1 Stelle ist und das kann eines von 4 möglichen Nukleosomen sein.

%\paragraph{F 11: Kommt H2A.Z bei allen Brustkrebsvarianten vor oder nur bei bestimmten?}

Weitere:
\paragraph{Kann man H2A.Z in macroH2A umwandeln?} Nicht beantwortet, Frage passt nicht zu Funktionalität.

\section{Genetic variation across and within individuals.}\label{sec:genetic-variation-across-and-within-individuals.}
% nicht mit dem Rücken zum Publikum, nicht einfach alles vorlesen aus Abbildungen...vor allem nicht einfach alles auf
% Deutsch vorlesen, was da auf eng steht, tabellen mit text sind keine guten Abbildungen, folienzahlen!
% Wtf mal nicht Leute unterbrechen, die Fragen stellen, rude
% Typen von Folgen (loss of function, etc.)

% Gut war Fazit & Rückbeziehung und Diskussionsfolie

% begriffe: extrinsische und intrinsische Motivationsstellen, somat. vs. Keimbahnmutationen (Art), typen mutationen:
% Substitutionen (SNV), Insert/Deletion < 50bp (Indels), ... weitere

% Art vs Typ bei Mutationen

Frage: Was ist SNV vs SNP

% große Abbildungen lieber auf Folien aufteilen statt zoomen

% Auswirkungen von Keimbahn- bzw. somat. Mutationen auf den Organismus

% man kann nicht immer sagen, ob insertion oder deletion vorliegt -  warum?


% In der Vorlesung haben wir gelernt, das unter bestimmtem Mutagenese-Einwirkungen (Bspw. UV vs andere) vorzugsweise
% in bestimmte Basen übergehen? Irgendwie sowas meinte Sagasser

\section{Plasmodesmata and intercellular molecular traffic control}\label{sec:plasmodesmata-and-intercellular-molecular-traffic-control}
Nüscht

\section{On the genetic basis of tail-loss evolution in humans and apes.}\label{sec:on-the-genetic-basis-of-tail-loss-evolution-in-humans-and-apes.}
Nüscht
% Vorwärtsgenetik vs. Rückwärtsgenetik



% Dritter Tag
\section{The emerging view on the origin and early evolution of eukaryotic cells}\label{sec:begin-of-eukaryotic-cells-oder-so}
% Zweidomänen-Modell
Nachschlagen: Protein-Struktur vergleichen statt Sequenz, um echte Neuheiten von unerkannten Homologen zu trennen
% "ich habe was ganz grundlegendes vergessen und frage einfach mal" come on

\section{Structure and Function of the 26S Proteasome}\label{sec:structure-and-function-of-the-26s-proteasome.}
% Proteasom, ubiquitin in defis packen, auch das Wort sterisch
RPN11 ist essenziell, hat man über Mutantenanalyse herausgefunden
